\begin{pastexam}[label={px:ncku-106-cce-ds-01}]{成大電通所 106 學年度 第一大題第 1 題}
（複選；5 分，全對才給分）

Which of the following statement(s) is (are) true?

(A)~The time complexity of a breath-first algorithm is O(V+E) which is the fastest search
algorithm.\\
(B)~The time complexity of Dijkstra algorithm is O(E lg V) if it is implemented by an array.\\
(C)~There may exist back edges in a spanning tree generated by a breath-first algorithm.\\
(D)~The time complexity to delete the maximum value in a heap is O(1).
\end{pastexam}

\begin{pastsoln}
答案：(C)。

原卷的 breath-first 是 breadth-first 的筆誤。 (C) 的「spanning tree 裡的 back edge」要讀成：
以 BFS 產生的生成樹（BFS 樹）為準，圖裡其他的邊可不可能是後向邊（指向樹上祖先的邊）。

\begin{itemize}
\item (A)：錯。 用鄰接串列時 BFS 是 $O(V + E)$（資結書第~21.3~節），這半句對；但「最快的搜尋
演算法」不成立。 DFS 也是 $O(V + E)$，而且任何要看過每個頂點和每條邊的走訪都至少要
$\Omega(V + E)$，沒有「BFS 最快」這回事；用鄰接矩陣時 BFS 還是 $\Theta(V^2)$。
\item (B)：錯。 $O(E \lg V)$ 是用二元堆積當優先佇列的時間（連通圖時 $O((V + E) \lg V) = O(E \lg V)$）。
用陣列時，每次找最小值要掃 $O(V)$、共 $V$ 次，鬆弛每次 $O(1)$，總共 $O(V^2 + E) = O(V^2)$
（資結書第~25.2~節）。
\item (C)：對。 無向圖的 BFS 不會有後向邊：若 $\{u, v\}$ 不是樹邊而 $u$ 是 $v$ 的祖先，$u$ 被取出時
就會發現 $v$（距離只差 $1$），$v$ 的父節點就會是 $u$，矛盾；所以無向圖 BFS 的非樹邊都是交叉邊。
但有向圖可以有：$1 \to 2$、$2 \to 1$ 從 $1$ 做 BFS，樹邊是 $1 \to 2$，$2 \to 1$ 指回祖先 $1$，是後向邊。
題目說「may exist」，有向圖的這個例子就讓它成立。
\item (D)：錯。 最大值在根，「找」是 $O(1)$；但「刪除」要把最後一個元素搬到根再往下調整
（sift-down），最壞走過整個高度，$O(\log n)$（資結書第~8.1~節）。
\end{itemize}
\end{pastsoln}
